\begin{helper}
Let \(V\) be a finite vertex set and fix \(v\in V\).  The coordinate-splitting
map
\[
q\longmapsto \bigl(q(v),(q(w))_{w\ne v}\bigr)
\]
from \(\mathbb R^V\) to \(\mathbb R\times\mathbb R^{V\setminus\{v\}}\)
preserves Lebesgue product measure.
\end{helper}
\begin{proof}
Choose the bijection \(\{\ast\}\sqcup (V\setminus\{v\})\to V\) sending the
distinguished point to \(v\) and every remaining coordinate to itself.  The
standard product-coordinate permutation induced by this bijection preserves
finite product Lebesgue measure.  The standard measurable equivalence
\(\mathbb R^{\{\ast\}\sqcup W}\cong \mathbb R^{\{\ast\}}\times\mathbb R^W\)
also preserves product measure, as does the identification
\(\mathbb R^{\{\ast\}}\cong\mathbb R\).  Their composition is exactly the
displayed coordinate-splitting map, so the map is measure preserving.
\end{proof}
